\documentclass{beamer}
\special{landscape}

%\usetheme{Berlin}
\usetheme{Warsaw}

%\usecolortheme{seahorse}
%\usefonttheme[onlysmall]{structurebold}

\setbeamertemplate{headline}[split]
\setbeamertemplate{footline}[default]
\setbeamertemplate{footline}[miniframes theme]
%\logo{\includegraphics[scale=0.25]{lifia_logo.png}}

\mode<presentation>
\usepackage[spanish]{babel}
\usepackage{beamerthemesplit}
\usepackage[utf8]{inputenc}
\usepackage{color}      % use if color is used in text


% Comandos en modo Verbatim
%\usepackage{fancyvrb}

\title{Práctica 3 - Parte 3 - USART}

\AtBeginSection[]

\begin{document}

\begin{frame}
\titlepage
\end{frame}

\section{USART - Introducción}
\subsection{Introducción}
\begin{frame}
\frametitle{USART}
\begin{itemize}
 \item Convertidor paralelo-serie y serie-paralelo.
 \item Permite a la CPU comunicarse con dispositivos serie
 \item Posee tres registros, de 8 bits.
  \begin{itemize}
   \item DIN: Registro de entrada, recibir datos de la línea serie.
   \item DOUT: Registro de salida, escribir datos a enviar por línea serie.
   \item CTRL: Registro de Control (escritura) y estado (lectura).
  \end{itemize}
 \item La recepción y envío de datos son independientes
 \item Los tres registros estan a partir de la posición 60h. 
  \begin{itemize}
      \item 60h = DIN
      \item 61h = DOUT
      \item 62h = CTRL
\end{itemize}
\end{itemize}
\end{frame}

\subsection{Interrupción}
\begin{frame}
\frametitle{USART - Interrupciones}
%\begin{block}{Slot Accounting System}
\begin{itemize}
 \item La USART puede generar 2 interrupciones
  \begin{itemize}
   \item INT 2: Cuando hay un carácter para ser recibido.
   \item INT 3: Cuando esta listo para enviar un carácter.
  \end{itemize}
  \item En la práctica no se va a considerar el uso de interrupciones
\end{itemize}
\end{frame}

\subsection{Caracteristicas de la Transmisión}
\begin{frame}
\frametitle{Caracteristicas}
%\begin{block}{Slot Accounting System}
\begin{itemize}
 \item Transmite 8bits por dato.
 \item Dos velocidades: V1: 6 baudios, V2: 18 baudios
 \item Comunicación Síncrona o asíncrona.
  \begin{itemize}
   \item Síncrona: 1 carácter de sincronismo, puede reconocer e insertar caracteres de sincronismo.
   \item Asíncrona: Sin paridad, 1 bit de parada y 1 bit de arranque.
  \end{itemize}
 \item La recepción y envío de datos son independientes
 \item Estas características se configuran en el registro de control
\end{itemize}
\end{frame}

\subsection{Registro de Conrol}
\begin{frame}
\frametitle{Registro de Control}
%\begin{block}{Slot Accounting System}
\begin{center}
 \includegraphics[scale=0.50]{usart-ctrl.png}
\end{center}

\begin{itemize}
 \item El registro de control 
 \begin{itemize}
   \item Sy/As: 0 = Síncrona / 1 = Asíncrona
   \item Vb: Velocidad: 0 = 6 baudios/ 1 = 18 baudios
   \item TxEN: Interrupción por TxRDY:  0 = inactivo / 1 = activo
   \item RxEN: Interrupción por RxRDY:  0 = inactivo / 1 = activo
   \item DTR: Data Terminal Ready, 0 inactivo, 1 activo
   \item RTS: Request to Send, 0 inactivo, 1 activo
   \item ER: Error Reset, 1 = Resetea flags de errores
   \item Sync: 1 = Inserción y búsqueda de caracteres de sincronismo (solo si Sy/AS = 0)
  \end{itemize}
\end{itemize}
\end{frame}

\begin{frame}
\frametitle{Registro de Estado}
%\begin{block}{Slot Accounting System}
\begin{center}
 \includegraphics[scale=0.50]{usart-estado.png}
\end{center}
\begin{itemize}
 \item El registro de estado 
 \begin{itemize}
   \item DSR: Indica estado de línea DSR
   \item SYNDET: Si Sy/AS = 0, y Sync = 1, indica recepción de carácter de sincronizacion
   \item CTS: Indica estado de entrada CTS
   \item FE: Indica error de recepción, bits de parada incorrectos
   \item OE: 1 = Se recibió un carácter y todavía no se habia leido el anterior 
   \item RxRDY: 1 = El registro de entrada tiene un dato
   \item TxRDY: 1 = El registro salida esta vacio, listo
   \end{itemize}
\end{itemize}
\end{frame}

\subsection{Líneas de Datos de USART}
\begin{frame}[fragile]
\frametitle{Impresora Serie}
\begin{itemize}
 \item USART posee esta lineas
 \begin{itemize}
  \item rxd: línea de salida, transmisión serie 
  \item txd: línea de entrada, recepción serie
  \item dsr: Data Set Ready, el dispositivo conectado esta listo para recibir/ enviar datos
  \item rts: Request to Send, la USART quiere enviar datos
  \item cts: Clear To Send, el dispositivo esta listo para recibir datos
 \end{itemize}
\end{itemize}
\end{frame}

\subsection{Conexión}
\begin{frame}[fragile]
\frametitle{Conexion USART - Impresora Serie}
\begin{center}
 \includegraphics[scale=1.00]{conexion_usart.png}
\end{center}

\end{frame}

\section{Impresora Serie}
\subsection{Introducción}
\begin{frame}[fragile]
\frametitle{Impresora Serie}
\begin{itemize}
 \item El simulador debe estar en la Configuración 4
 \item Utiliza una interfaz serie RS-232 simplificada
 \item Puede recibir hasta 5 caracteres (buffer)
 \item Caracteristicas de Comunicación
 \begin{itemize}
  \item Comunicacion Asincronica
  \item Carácteres de 8 bits
  \item Dos velocidades (6 baudios y 18 baudios)
  \item Sin paridad
  \item Un bit de arranque (0)
  \item Un bit de parada (1)
 \end{itemize}
  \item Dos protocolos de Comunicación
\begin{itemize}
 \item Hardware DTR
 \item Software XON/XOFF
\end{itemize}
\item Cambiamos la velocidad de la impresora con el comando ''vi''
\end{itemize}
\end{frame}

\begin{frame}[fragile]
\frametitle{Conexion USART - Impresora Serie}
\begin{itemize}
 \item La conexion USART e Impresora usa las siguientes lineas:
 \begin{itemize}
  \item TxD: línea de Transmisión serie, la CPU transmite los datos a imprimir
  \item RxD: línea de Recepción serie, los datos enviados desde la impresora a la CPU
  \item DSR: línea de salida que habilita/inhabilita la transmisión de datos a la impresora
\end{itemize}
\item Las lineas de RTS y CTS de la USART estan conectadas entre si (no interviene la impresora)
\end{itemize}
\end{frame}

\section{Protocolos de Comunicación}
\subsection{Hardware DTR}
\begin{frame}[fragile]
\frametitle{DTR}
\begin{itemize}
 \item Configurar la impresora para usar el protocolo DTR usando el comado `` pi d'' 
 \item La línea DTR indica si se debe iniciar/suspender la transmisión de datos.
 \item Si el buffer esta lleno, la impresora pone 0 en la línea DTR.
 \item Si el buffer tiene lugar, la impresora pone 1 en la línea DTR.
 \item Puede ser mediante Interrupciones o consulta de estado
 \item La línea DTR de la impresora está conectada a la linea DSR de la USART
\end{itemize}

\end{frame}

\subsection{Software XON/XOFF}
\begin{frame}[fragile]
\frametitle{XON/XOFF}
\begin{itemize}
 \item Configurar la impresora para usar el protocolo DTR usando el comado `` pi x''
 \item Utiliza unicamente TxD y RxD.
 \item Cuando el buffer de recepeción esta lleno, la impresora envía el carácter XOFF a través de la línea TxD, para detener la transmisión de datos
 \item Cuando hay espacio en el buffer, la impresora envía el carácter XON, avisando que se puede seguir enviando caracteres
 \item El Carácter XOFF es 13h
 \item El Carácter XON es 11h
 \item Puede ser mediante Interrupciones o consulta de estado
\end{itemize}

\end{frame}



\section{Protocolo DTR por Consulta de Estado}

\subsection{Configuración}
\begin{frame}[fragile]
\frametitle{Como configurar la Impresora con DTR}
\begin{itemize}
 \item El simulador debe estar en la configuración 4
 \item Se configura la USART con el numero binario 01010001
 \begin{itemize}
 \item Error reset (bit 6 en 1)
 \item DTR activado (bit 4 en 1)
 \item Comunicación asíncrona (bit 0 en 1)
 \item Velocidad: 6 baudios (bit 1 = 0).
\end{itemize}
\item Utilizamos velocidad de 6 baudios, si cambiamos la velocidad, hay que cambiar la velocidad de la impresora con el comando ''vi''
\end{itemize}
\end{frame}

\begin{frame}[fragile]
\frametitle{Ejemplo}
\begin{block}{Configuración}
 \begin{verbatim}
         CTRL EQU 62H
         
           ...
           
         MOV AL, 01010001B ; = 51H
         OUT CTRL, AL
 \end{verbatim}
\end{block}

\end{frame}

\subsection{Impresión}

\begin{frame}[fragile]
\frametitle{Como usar la Impresora con DTR}
\begin{itemize}
 \item  Consultamos para saber si esta lista para recibir un carácter. Se deben dar estas condiciones:
 \begin{itemize}
  \item Data Set Ready = 1, la impresora debe estar lista para recibir carácter (bit de estado 7)
  \item DSR está conectada a la línea DTR de la impresora
  \item TxRDY = 1, la USART debe estar lista para transmitir (bit de estado 0) 
  \item TxRDY vale 1 si CTS vale 1 y no se esta trasmitiendo ya algun dato
 \end{itemize}
\item Transmitimos el siguiente carácter 
\item Volvemos a empezar si nos quedan mas caracteres a enviar a la impresora.
\end{itemize}

\end{frame}


\begin{frame}[fragile]
\frametitle{Ejemplo}
\begin{block}{Esperar a TxRDY y DSR en 1}
 \begin{verbatim}
         CTRL EQU 62H
         
           ...
           
TEST: IN AL, CTRL
      AND AL, 10000001B ; = 81H
      CMP AL, 81H ; Porque hacemos CMP?
      JNZ TEST
      
           ...
 \end{verbatim}
\end{block}

\end{frame}

\begin{frame}[fragile]
\frametitle{Envio del carácter}
\begin{block}{Enviar el carácter}
 \begin{verbatim}
        DOUT EQU 61H
         
        ...
           
       MOV AL, [BX] ; Proximo carácter
       OUT DOUT, AL
             
 \end{verbatim}
\end{block}

\end{frame}


\begin{frame}[fragile]
\frametitle{Ejercicio}

\begin{itemize}
 \item Veamos el Ejercicio 4b de la Práctica
\end{itemize}

\end{frame}

\section{Protocolo XON/XOFF por Consulta de Estado}

\subsection{Configuración}

\begin{frame}[fragile]
\frametitle{Como configurar la Impresora con XON/XOFF}
\begin{itemize}
 \item El simulador debe estar en la configuración 4
 \item Se configura la USART con el numero 51H (en binario 01010001 en el registro de control.
 \begin{itemize}
 \item Error reset activo (bit 6)
 \item DTR activado (bit 4)
 \item Comunicación asíncrona (bit 0). 
 \item Velocidad: 6 baudios (bit 1 = 0).
\end{itemize}
\item Utilizamos velocidad de 6 baudios, si cambiamos la velocidad, hay que cambiar la velocidad de la impresora con el comando ''vi''
 \end{itemize}
\end{frame}

\begin{frame}[fragile]
\frametitle{Ejemplo}
\begin{block}{Configuración}
 \begin{verbatim}
         CTRL EQU 62H
         
           ...
           
         MOV AL, 01010001B ; = 51H
         OUT CTRL, AL
 \end{verbatim}
\end{block}

\end{frame}

\subsection{Impresión}
\begin{frame}[fragile]
\frametitle{Como usar la Impresora con XON/XOFF}

\begin{itemize}
\item Esperamos a que la USART este lista para enviar un carácter. 
\begin{itemize}
\item Es decir TxRDY cambia a 1 
\end{itemize}
\item Enviamos el carácter a imprimir. 
\item Verificar si recibimos un carácter
\begin{itemize}
 \item Es decir RxRDY cambia a 1
\end{itemize}
\item Si no recibimos un carácter, podemos enviar el siguiente carácter
\item Si recibimos un carácter, puede ser XON (11h) o XOFF (13h)
\begin{itemize}
\item Si recibimos XON, podemos seguir imprimiendo.
\item Si recibimos XOFF, hay que esperar hasta recibir XON para poder imprimir.
\end{itemize}
\end{itemize}

\end{frame}

\begin{frame}[fragile]
\frametitle{Esperar que este lista la impresora}
\begin{block}{Esperar Impresora}
 \begin{verbatim}
        CTRL EQU 62H
         
        ...
           
TEST:  IN AL, CTRL
       AND AL, 1
       JZ TEST
             
 \end{verbatim}
\end{block}

\end{frame}

\begin{frame}[fragile]
\frametitle{Envio del carácter}
\begin{block}{Enviar el carácter}
 \begin{verbatim}
        DOUT EQU 61H
         
        ...
           
       MOV AL, [BX] ; Proximo carácter
       OUT DOUT, AL
             
 \end{verbatim}
\end{block}

\end{frame}

\begin{frame}[fragile]
\frametitle{Recibimos carácter?}
\begin{block}{Hay un dato disponible?}
 \begin{verbatim}
        CTRL EQU 62H
         
        ...
        IN AL, CTRL
        AND AL, 02H     : Bit 1 RxRDY
        JZ HAY CARACTER
             
 \end{verbatim}
\end{block}

\end{frame}

\begin{frame}[fragile]
\frametitle{Leer carácter}
\begin{block}{Leer el dato disponible}
 \begin{verbatim}
        DIN EQU 60H
         
        ...
        IN AL, DIN ; AL tiene el carácter recibido
        
 \end{verbatim}
\end{block}

\end{frame}

\begin{frame}[fragile]
\frametitle{Ejercicio}

\begin{itemize}
 \item Veamos el Ejercicio 4c de la Práctica
\end{itemize}

\end{frame}

\end{document}
